\section{Alternating fractional memories}
\label{sec:alternating}

The second mechanism alternates between two coordinates: the current
configuration is read from one coordinate, which stays fixed while the
next configuration is appended to the other.
Neither coordinate is erased.  The available precision increases with
the slot index; the displacement needed to write the next block is much
smaller than every fixed derivative cost of reading the current block.

\begin{theorem}[Rapid decay with alternating memories]
\label{thm:alternating}
Fix a positive computable viscosity. From a deterministic Turing machine
and a finite input one can effectively
construct a smooth force $f$ on $[0,\infty)\times\mathbb R^3$, with
spatial support contained in a fixed compact set independent of time,
such that for every $J,l\geq0$ and spatial multi-index $\alpha$,
\begin{equation}
 \sup_{t\geq0,\,X\in\mathbb R^3}
 (1+t+|X|)^J
 |\partial_t^l\partial_X^\alpha f(t,X)|<\infty.
 \label{eq:alternating-schwartz}
\end{equation}
The solution with zero initial velocity is explicit, global, and unique
in the class \eqref{eq:energy-class}, with pressure zero.  Its velocity satisfies
the same support and derivative estimates.  The material particle
initially at $(-1,0,0)$ reaches $X_1>0$ if and only if the machine halts.
\end{theorem}

The proof has three tasks: recover one complete configuration from a
stationary donor coordinate, use it to prescribe the next displacement,
and bound the derivatives of that displacement.

\subsection{A finite block and an alternating history}
Let $Q$ be the machine's state set, $H_Q\subseteq Q$ its halting states,
and $\Gamma$ its tape alphabet. Normalize missing transitions by adding
a halting state, without moving or changing the scanned symbol.
Extend the resulting table after a halting state by
$\delta(q,a)=(q,a,0)$.  This extension merely freezes a halted
configuration.  Let $N$ be the input length, choose
$b=2(1+\max(|Q|,|\Gamma|))$, and assign distinct even digits
$g_a,g_q\in\{0,2,\ldots,b-2\}$ within the symbol and state families,
with blank digit zero.  Write
$m=(b-2)/(b-1)<1$ and
$J_s=[g_s/b,(g_s+1)/b]$ for each member of either family.
Intervals are disjoint within each family.  Using rational rescalings
of the effective cutoffs in Section~\ref{sec:analytic}, choose smooth functions
$\phi_s=1$ on a neighborhood of $J_s$, with pairwise disjoint supports
within that family and effective bounds for every derivative.

We work in coordinates $(x,y,z)=(X_2,X_3,X_1)$, a right-handed cyclic
permutation.  After $n$ machine transitions let $\xi_n,\eta_n$ be the
fractional base-$b$ codes of the tape strictly left of the head and from
the head rightward, each nearest symbol first.  They need at most
$K_n=N+2n+2$ digits, padded by zeros.  Indeed the head moves at most
$n$ cells, while a nonblank cell was either in the initial input or was
visited in those $n$ moves.  Put
\begin{equation}
 C_n=\frac{g_{q_n}+\xi_n+b^{-K_n}\eta_n}{b},\qquad
 s_n=(N+1)2^{(n+3)^2},\qquad \epsilon_n=b^{-s_n}.
 \label{eq:alternating-block}
\end{equation}
The block has one state digit, $K_n$ left-stack places, and $K_n$
right-stack places, with zero padding. Every digit is allowed, so all
fractional tails lie in $[0,m]$, and
$b^{1+2K_n}C_n$ is an integer.  We have
\[
 s_n>1+2K_n,\qquad
 \frac{s_{n+1}}{s_n}=2^{2n+7},\qquad
 s_{n+1}>s_n+1+2K_n.
\]
For example $2^{(n+3)^2}\geq8(n+1)$ gives the first inequality directly.
The intended memories at time $1+n$ are
\begin{equation}
 x_n=\sum_{\substack{0\leq j\leq n\\j\text{ even}}}\epsilon_jC_j,
 \qquad
 y_n=\sum_{\substack{0\leq j\leq n\\j\text{ odd}}}\epsilon_jC_j.
 \label{eq:alternating-history}
\end{equation}
These are specifications of the trajectory to be proved, not parameters
obtained by executing the machine while constructing the field.
Figure~\ref{fig:alternating} illustrates the stationary-donor invariant.

\begin{figure}[ht]
\centering
\begin{tikzpicture}[>=Latex,font=\small]
\node[draw,rounded corners,align=center,minimum width=3.2cm,minimum height=1.1cm]
 (a) at (0,0) {$x_n$ contains $C_n$\\$y_n$ contains earlier records};
\node[draw,rounded corners,align=center,minimum width=3.2cm,minimum height=1.1cm]
 (b) at (5.2,0) {$x_n$ unchanged\\$y_n+\epsilon_{n+1}C_{n+1}$};
\node[draw,rounded corners,align=center,minimum width=3.2cm,minimum height=1.1cm]
 (c) at (5.2,-2) {$x_n+\epsilon_{n+2}C_{n+2}$\\$y_{n+1}$ unchanged};
\draw[->,thick] (a)--node[above=8mm]{read $x$; write $y$}(b);
\draw[->,thick] (b)--node[right]{read $y$; write $x$}(c);
\end{tikzpicture}
\caption{Two successive slots, starting with even $n$. Every old record
remains present. The donor does not move during its slot, so its decoded
instruction stays constant while the new record is written. The signal
coordinate $z$ is omitted from the diagram.}
\label{fig:alternating}
\end{figure}

\subsection{A smooth decoder on every real donor value}
The positional tape representation is a generalized-shift encoding
\cite{Moore}. We now give the decoder on arbitrary real arguments,
so the field is defined independently of the distinguished execution.
Choose rational $m<a_P<b_P<1$ and define on $[0,1]$
\[
 P(v)=v-H\bigl((v-a_P)/(b_P-a_P)\bigr).
\]
Near zero this is $v$, and near one it is $v-1$; it therefore extends
to a smooth bounded $1$-periodic function, with all derivatives bounded.
It equals the fractional part whenever that fractional part belongs to
$[0,m]$.  For slot $n$, let the donor $D$ be $x$ when $n$ is even and
$y$ when $n$ is odd, and set
\begin{equation}
 C=P(b^{s_n}D),\qquad B=P(b^{s_n+1+K_n}D),\qquad
 A=P(b^{s_n+1}D)-b^{-K_n}B.
 \label{eq:alternating-read}
\end{equation}
At the intended donor these give $C_n,\eta_n,\xi_n$ respectively.
To verify this, every earlier donor block with index $j<n$ contributes
an integer after multiplication by $b^{s_n}$: its denominator is
$b^{s_j+1+2K_j}$, and
$s_n\geq s_{j+1}>s_j+1+2K_j$.  The remaining block is $C_n$.
Multiplication once more by $b$ removes the state digit.  Multiplication
by $b^{K_n}$ removes the first $K_n$ tape digits since
$b^{K_n}\xi_n$ is an integer, leaving the tail $\eta_n$.
All fractional parts used here lie in $[0,m]$, where $P$ is exact.
Subtraction then gives $A=\xi_n$.  This is the precise role of old-scale
divisibility in the alternating encoding.

Make a finite list of the normalized table branches, including its
freezing halting transitions, and split each left move by its left-top
symbol $\ell$. For a branch
$i:(q,a)\mapsto(p_i,b_i,d_i)$ put
\[
 \Psi_i(D)=\phi_q(C)\phi_a(B)
 \begin{cases}1,&d_i=0,1,\\\phi_\ell(A),&d_i=-1.\end{cases}
\]
Let $R=bB-g_a$ and define its tape update by
\begin{equation}
 (T_i^1,T_i^2)=
 \begin{cases}
 ((g_{b_i}+A)/b,R),&d_i=1,\\
 (A,(g_{b_i}+R)/b),&d_i=0,\\
 (bA-g_\ell,(g_\ell+(g_{b_i}+R)/b)/b),&d_i=-1.
 \end{cases}
 \label{eq:alternating-update}
\end{equation}
At an intended donor precisely the correct branch has weight one, and
all others vanish.  The displayed digit push and pop formulas give
$\xi_{n+1},\eta_{n+1}$, which fit in $K_{n+1}$ places.  Define, on
all real donor arguments,
\begin{align}
 L_n(D)&=\frac{\epsilon_{n+1}}b
    \sum_i\Psi_i(D)
       \bigl(g_{p_i}+T_i^1+b^{-K_{n+1}}T_i^2\bigr),
       \label{eq:alternating-write}\\
 E_n(D)&=\epsilon_n-\epsilon_{n+1}
       +2\epsilon_{n+1}\sum_{q\in H_Q}\phi_q(C).
       \label{eq:alternating-signal}
\end{align}
Thus $L_n(D)=\epsilon_{n+1}C_{n+1}$ on the intended donor.
The field retains all old records, so the branches need no disjoint
images and no reversible-table preprocessing.

\subsection{Keeping the donor stationary during the write}
Choose one fixed compact cutoff $\chi$ equal to one on a neighborhood
of $[0,1]^2\times[-1,1]$ in $(x,y,z)$ coordinates, again using rational
rescalings of the effective cutoffs so that every derivative has an
effective bound. In slot
$[1+n,2+n]$ prescribe $U=\theta'(t-1-n)V_n$, where
\begin{equation}
 V_n=
 \begin{cases}
 \operatorname{curl}(\chi(zL_n(x)-yE_n(x)),0,0),&n\text{ even},\\
 \operatorname{curl}(0,\chi(xE_n(y)-zL_n(y)),0),&n\text{ odd}.
 \end{cases}
 \label{eq:alternating-curls}
\end{equation}
On the plateau these velocities are respectively $(0,L_n,E_n)$ and
$(L_n,0,E_n)$.  In particular the donor is stationary throughout its
slot. During initialization, for $0\leq t\leq1$, set
\[
 d=(\epsilon_0C_0,0,1-\epsilon_0),\qquad
 U(t,x,y,z)=\theta'(t)\operatorname{curl}
 \left[\frac12\chi(x,y,z)\bigl(d\times(x,y,z)\bigr)\right].
\]
On the plateau this equals $\theta'(t)d$, so the particle follows
$(0,0,-1)+\theta(t)d$. This whole segment lies in the plateau.
Only the initial code $C_0$ is supplied to the field.

We verify the trajectory before using it for observation.  Initialization
ends at $(\epsilon_0C_0,0,-\epsilon_0)$.  If the memories at the
start of slot $n$ are \eqref{eq:alternating-history}, the
stationary donor gives the values in \eqref{eq:alternating-read}.
With progress $\theta(t-1-n)$, the other memory acquires $L_n(D)$
and the signal acquires $E_n(D)$.  This proves the next memory identity.
All proposed paths remain on the plateau: the two memories are between
zero and $\sum_j\epsilon_j$, and
\[
 \sum_{j\geq0}\epsilon_j\leq2\epsilon_0<1/4.
\]
Here $s_{j+1}-s_j\geq1$ and $b\geq4$ already give a geometric
majorant; the stated constant two is more than sufficient.  The signal
starts at $-\epsilon_0$, and each increment lies between zero and
$\epsilon_n+\epsilon_{n+1}$.  Its total displacement is less than
$2\sum_j\epsilon_j<1/2$.  Hence $|z|<1$ on all slots, even after
halting.  The proposed curves genuinely solve the cutoff field, and
ODE uniqueness identifies them with the material trajectory.

Before the first slot whose current state is halting, the signal is
$-\epsilon_n$ at the slot's start and interpolates monotonically to
$-\epsilon_{n+1}$, remaining strictly negative.  In the first halting
slot it instead ends at $+\epsilon_{n+1}$.  This also covers an initially
halting machine, which is detected in the first post-loading slot.
Loading stays negative.  Thus $z=X_1>0$ occurs exactly when the machine
halts, with no intermediate false positive.

\subsection{Derivative bounds and effective evaluation}
The particle correspondence has been proved for every real time. It
remains to prove that the field defining it has the claimed rapid bounds.  The
functions $C,A,B$ in \eqref{eq:alternating-read} are uniformly
bounded.  Applying the chain and product rules through spatial order
$r$ gives
\[
 \left\|L_n/\epsilon_{n+1}\right\|_{C^r}
 +\left\|\sum_{q\in H_Q}\phi_q(C)\right\|_{C^r}
 \leq C_r b^{r(s_n+1+K_n)}.
\]
Constants depend on the fixed table, base, and smooth functions, but not
on $n$.  The factors $b^{-K_n}$ and $b^{-K_{n+1}}$ are at most one;
there are only finitely many branches.  Accounting for one derivative
in the curl and for the fixed cutoff yields
\begin{equation}
 \|V_n\|_{C^r}
 \leq C_r'\left(b^{-s_n}
      +b^{-s_{n+1}+(r+1)(s_n+1+K_n)}\right).
 \label{eq:alternating-estimate}
\end{equation}
For each fixed $r$, the second term is at most $b^{-s_n}$ for all
sufficiently large $n$, because $s_{n+1}/s_n=2^{2n+7}$ and
$K_n+1<s_n$.  Temporal derivatives only differentiate the fixed
clock pulse, and the zero collars give smooth gluing.  The right side
therefore tends to zero faster than every inverse power of $1+n$ for
each fixed mixed order.  Early slots and loading have finite individual
bounds.  These facts prove all the rapid mixed bounds for $U$.

For $f=\mathcal R_\nu[U]$, the linear terms require one additional
time derivative or two additional spatial derivatives.  Each derivative
of the quadratic term is a finite sum of products of mixed derivatives
of $U$, one with an additional spatial derivative.  The already proved
estimates bound every such factor; one factor may carry any requested
temporal weight and the other is bounded.  Consequently $f$ has every
rapid mixed bound.  Fixed compact support turns a temporal weight into
the full weight $(1+t+|X|)^J$ in
\eqref{eq:alternating-schwartz}.  The fixed support gives all Sobolev conditions in
\eqref{eq:energy-class}. Lemma~\ref{lem:realization} supplies
uniqueness and the global material flow.

For a finite slot index all scales, the fixed branch list, and all
coefficients are effective.  A coarse enclosing time interval selects
finitely many candidate slots; zero collars make evaluation consistent
where such choices overlap.  The periodic function $P$ can be evaluated
by finitely many neighboring integer translates, or by its affine
formula across an integer, without a discontinuous fractional-part
test.  Flat cutoff bounds and finite differentiation give arbitrary
rational-error evaluation of the field, the force, and every derivative.
Bounds for any fixed finite set of orders are obtained from a finite
initial range plus the explicit tail inequality in
\eqref{eq:alternating-estimate}.  Only $C_0$ is computed from
the input orbit; subsequent slot functions perform one local table
update on their spatial donor arguments.  Physical slots have unit
length and no finite accumulation time. This proves
Theorem~\ref{thm:alternating}.
